\section{Vật lí, môn khoa học cơ bản nhất}

\ 

Vật lí không chỉ là một ngành khoa học bất kì, mà còn là ngành khoa học tự nhiên cơ bản nhất. Do các đối tượng được nghiên cứu của những ngành khác, bao gồm chuyển động, lực, công suất, năng lượng, nhiệt, âm thanh, ánh sáng, từ thế giới vi mô đến vĩ mô, cũng đều được quan tâm bởi các nhà vật lí học. Lấy một ví dụ điển hình, sinh học thường xuyên nghiên cứu về quá trình chuyển hóa chất, mà khi đã nói về sự tương tác giữa các phân tử và nguyên tử thì chắc chắn phải đề cập đến hóa học. Mặc dù vậy, khi nhìn sâu hơn, tuy hóa học cho chúng ta hiểu về các phản ứng giữa các nguyên tử, để có thể thấy được lí do sâu sa dẫn đến phản ứng và quan sát được sự tương tác của bộ phận cấu tạo những nguyên tử đó, chúng ta sẽ cần những kiến thức vật lí về nguyên tử và hạt nhân. Đương nhiên, chúng ta sẽ không kì vọng sử dụng các định luật vật lí cơ bản để giải thích trực tiếp các quá trình phức tập của hóa và sinh, giống như chúng ta không phân tích một cuốn sách bằng việc phân tích vị trí của từng chữ cái, nhưng lượng kiến thức chúng ta sẽ tích cóp được sẽ được sử dụng để hiểu hơn về các cấu trúc trong khoa học tự nhiên sau này.